\DocumentMetadata{
  pdfversion=2.0,
  pdfstandard=ua-2,
  lang=en-US,
  tagging=on,
  tagging-setup={math/setup=mathml-SE,tabsorder=structure}
}
\documentclass[11pt,letterpaper]{article}
\usepackage[margin=0.68in,includefoot]{geometry}
\usepackage{newcomputermodern}
\usepackage{amsmath,amscd,mathtools}
\usepackage{tikz,adjustbox}
\providecommand{\square}{\mdlgwhtsquare}
\usepackage[hidelinks]{hyperref}
\hypersetup{
  pdftitle={Combinatorics PhD Exam, January 2002},
  pdfauthor={University of Florida Department of Mathematics},
  pdfsubject={Graduate mathematics examination},
  pdfdisplaydoctitle=true,
  bookmarksopen=true
}
\setcounter{secnumdepth}{0}
\setlength{\parindent}{0pt}
\setlength{\parskip}{0.28em}
\setlength{\emergencystretch}{3em}
\setlength{\leftmargini}{2.15em}
\setlength{\leftmarginii}{2.65em}
\setlength{\leftmarginiii}{3em}
\renewcommand{\labelenumi}{\textbf{\theenumi.}}
\pagestyle{empty}
\raggedbottom
\allowdisplaybreaks
\newenvironment{examproblems}[1]{%
  \begin{enumerate}%
  \setcounter{enumi}{#1}%
  \setlength{\topsep}{0.35em}%
  \setlength{\partopsep}{0pt}%
  \setlength{\itemsep}{0.48em}%
  \setlength{\parsep}{0.18em}%
}{\end{enumerate}}
\newenvironment{examsubparts}{%
  \begin{enumerate}%
  \setlength{\topsep}{0.18em}%
  \setlength{\partopsep}{0pt}%
  \setlength{\itemsep}{0.16em}%
  \setlength{\parsep}{0.1em}%
}{\end{enumerate}}
\newenvironment{examinstructions}{%
  \par\begingroup\itshape
}{%
  \par\endgroup\vspace{0.18em}
}
\makeatletter
\renewcommand\section{\@startsection{section}{1}{0pt}%
  {-1.25ex plus -.25ex minus -.1ex}%
  {0.55ex plus .1ex}%
  {\normalfont\large\bfseries}}
\renewcommand{\@maketitle}{%
  \newpage\null\vskip -1em%
  {\footnotesize\bfseries
    \href{https://www.ufl.edu/}{University of Florida}, \href{https://math.ufl.edu/}{Department of Mathematics}\par}%
  \vskip -0.5em%
  \tagstructbegin{tag=Title}%
    {\LARGE\bfseries\href{https://gma.math.ufl.edu/exams/combinatorics-phd/}{Combinatorics PhD Exam}, January 2002\par}%
  \tagstructend%
  \vskip 0.25em%
  \tagmcbegin{artifact}%
    \hrule height 1pt%
  \tagmcend%
  \vskip 1em%
}
\makeatother
\title{Combinatorics PhD Exam, January 2002}
\author{}
\date{}
\begin{document}
\maketitle
\thispagestyle{empty}
\begin{examinstructions}
SHOW ALL WORK TO RECEIVE CREDIT!!
\end{examinstructions}
\begin{examproblems}{0}
\item
Find an explicit formula for \(a_k\) if \(a_0=1\) and \(a_{k+1}=3a_k+2\).
\item
Let \(a_0=1\). Find the unique sequence of real numbers satisfying 
\[
\sum_{i=0}^{n} a_i a_{n-i}=1
\]
 for all natural numbers~\(n\).
\item
A dealership has~\(n\) cars. An employee with a sense of humor takes all~\(n\) keys, puts one of them in each car at random, then locks the doors of all cars. When the owner of the dealership discovers the problem, he calls a locksmith. He tells him to break in to a car, then use the key found in that car to open another, and so on. If and when the keys already recovered by this procedure cannot open any new cars, the locksmith is to break in to another car. This algorithm goes on until all cars are open. What is the probability that the locksmith will only have to break in to one car?
\item
We have~\(n\) distinct cards. We want to split their set into nonempty subsets so that each of them contains an even number of cards. Then we want to order the cards within each subgroup. Finally, we want to order these subgroups into a line. Find an explicit formula for the number of ways \(g_n\) we can do this.
\item
How many permutations of length~\(8\) are there with descent set \(\{3,6\}\)?
\item
Let \(P_n\) be the poset of all parking functions of size~\(n\), with the coordinate-wise ordering. That is, \(x\leq y\) if \(x(i)\leq y(i)\) for all~\(i\). Add a maximum element \(\hat{1}\) to the top of \(P_n\) to get the poset \(P'\). Is \(P'\) a lattice?
\item
Is it true that a finite poset has as many ideals as antichains?
\item
Let \(f_k(n)\) be the number of permutations of length~\(n\) having~\(k\) descents. Prove that for any fixed~\(k\), \(f_k(n)\) is a~\(p\)-recursive function of~\(n\).
\item
Four friends, \(A\), \(B\), \(C\), and~\(D\) organize a long jump competition every day until the final order of the four of them will be the same on two different days. At most how long will they have to wait for that to happen? (Each jump is measured in centimeters, so all kinds of ties, twofold, threefold, fourfold, are possible).
\item
Prove that for any positive integer~\(n\), we have 
\[
\mu_{NC_n}(\hat{0},\hat{1})=(-1)^{n-1}c_n.
\]
 Here \(c_n\) is the~\(n\)th Catalan number, and \(NC_n\) is the lattice of all noncrossing partitions of~\(n\).
\end{examproblems}
\end{document}
